\chapter{Their sets, defined, and what runs on them exactly}

\textbf{Purpose:} Define the sets Fefferman's problem statement names, as sets; record what he says
was known and what he asks for; and record what the engine's exact integer arithmetic shows when it
is pointed at those sets, with the boundary of each object read off the object by a probe and the
inheritance of that boundary through the constructors measured. \textbf{Scope:} Fefferman's
statement, six pages with the errata, fetched from the Clay site and read in full on 2026-09-17;
\texttt{examples/\allowbreak{}0\_experimental/\allowbreak{}exact\_navier\_stokes\_on\_torus.py}
and \texttt{evidence/\allowbreak{}proofs/\allowbreak{}posits/\allowbreak{}proof\_boundary\_inheritance.py}
in the \texttt{anchor\_sift} tree, both run and both exiting 0; and the dated workbook beside them at
\texttt{theory/\allowbreak{}workbook/\allowbreak{}navier\_stokes\_workbook.md}, whose withdrawn
entries this chapter repeats.

Nothing in this chapter bears on whether (A), (B), (C) or (D) holds. It follows the rail the first
chapter set down: claim nothing, cite nothing unread, put the gap in the sentence that makes the
claim, keep a withdrawn claim on the page with what killed it, draw the bar and never derive it. It
adds one more, from the arithmetic: the boundary kind of an object is not assigned by judgment. It is
read off the object by a probe, and the object answers.

We know nothing here that the field does not. Every object in this chapter belongs to someone else
and is named with respect in the prior-art section. The two files run those objects in exact
arithmetic, which adds no rounding and a second route, and they add nothing else.

\section{Their sets, defined}

Written as sets so that every later sentence names the set it is about. $n = 3$ and $\nu > 0$ are
fixed throughout, $e_j$ are the unit vectors, and smooth means $C^\infty$. Each set is Fefferman's
condition as he wrote it, and the definitions are transcribed, not designed. His numbered conditions
are: (1) the momentum equation, (2) $\operatorname{div} u = 0$, (3) $u(x,0) = u^\circ(x)$, (4)
rapid decay of every derivative of the datum on $\R^3$, (5) rapid decay of every space-time
derivative of the force on $\R^3\times[0,\infty)$, (6) $p, u$ smooth there, (7) one constant
bounding $\int_{\R^3}|u(x,t)|^2\,dx$ for every $t\ge 0$, (8) the datum and the force periodic in each
$e_j$, with the errata adding that the pressure is periodic, (9) rapid decay of every derivative of
the force in $t$ alone, (10) the solution periodic, (11) $p, u$ smooth.

\begin{itemize}
\item $D_4$, the whole-space data: smooth $u^\circ:\R^3\to\R^3$ with $\operatorname{div}u^\circ = 0$
  and, for every multi-index $\alpha$ and every $K>0$, a constant $C_{\alpha K}$ with
  $|\partial_x^\alpha u^\circ(x)| \le C_{\alpha K}(1+|x|)^{-K}$ on $\R^3$.
\item $F_5$, the whole-space forces: smooth $f:\R^3\times[0,\infty)\to\R^3$ with, for every
  $\alpha, m, K$, a constant with $|\partial_x^\alpha\partial_t^m f(x,t)| \le C_{\alpha m K}(1+|x|+t)^{-K}$.
\item $S_{67}(u^\circ,f)$, the physically reasonable whole-space solutions for a datum and a force:
  pairs $(p,u)$ smooth on $\R^3\times[0,\infty)$, satisfying (1) and (2) there, with $u(x,0)=u^\circ(x)$
  and one constant $C$ with $\int_{\R^3}|u(x,t)|^2\,dx < C$ for every $t \ge 0$.
\item $D_8$, the periodic data: smooth $u^\circ$ with $\operatorname{div}u^\circ=0$ and
  $u^\circ(x+e_j) = u^\circ(x)$ for $j = 1,2,3$. No decay condition; the torus replaces it.
\item $F_{89}$, the periodic forces: smooth $f$ with $f(x+e_j,t) = f(x,t)$ and, for every
  $\alpha, m, K$, a constant with $|\partial_x^\alpha\partial_t^m f(x,t)| \le C_{\alpha m K}(1+|t|)^{-K}$.
\item $S_{1011}(u^\circ,f)$, the periodic solutions: pairs $(p,u)$ smooth on $\R^3\times[0,\infty)$,
  satisfying (1), (2), (3), with $u(x+e_j,t) = u(x,t)$ and, by the errata, $p(x+e_j,t) = p(x,t)$. No
  energy condition is written; on a period cell the energy of a smooth periodic field is finite.
\end{itemize}

The four alternatives in those names:

\begin{itemize}
\item (A): for every $u^\circ$ in $D_4$, $S_{67}(u^\circ, 0)$ is not empty.
\item (B): for every $u^\circ$ in $D_8$, $S_{1011}(u^\circ, 0)$ is not empty.
\item (C): there exist $u^\circ$ in $D_4$ and $f$ in $F_5$ with $S_{67}(u^\circ,f)$ empty.
\item (D): there exist $u^\circ$ in $D_8$ and $f$ in $F_{89}$ with $S_{1011}(u^\circ,f)$ empty.
\end{itemize}

Two readings of those definitions, both from the statement and neither new. (C) is not the negation
of (A): the negation of (A) is (C) with $f = 0$, and (C) allows any $f$ in $F_5$, so (A) and (C) can
both be true, and likewise (B) and (D). The previous chapter called that two different removals, and
the removal instances below are that shape on one field each. Fefferman's own sentence on the
unforced case is that either (A) and (B) hold, or else there is a smooth divergence-free $u^\circ$
for which (1), (2), (3) have a solution with a finite blowup time $T$, past which the velocity becomes
unbounded; for Euler, $\nu = 0$, the Beale--Kato--Majda criterion then gives
$\int_0^T \sup_x |\operatorname{curl} u(x,t)|\,dt = \infty$. Both are his sentences and nothing here
touches them.

Two more sets he defines and this chapter does not use: the weak solutions, pairs $(p,u)$ with $u$
in $L^2$, $p$ and $f$ in $L^1$, satisfying his (12) for every smooth compactly supported test field
and (13) for every smooth compactly supported test function, which Leray showed always exist in three
dimensions with suitable growth; and the singular set of a weak solution, the points near which $u$ is
unbounded, whose one-dimensional parabolic Hausdorff measure Caffarelli, Kohn and Nirenberg showed
is zero, with F.-H. Lin's simpler proof. Named here because they are in the statement, and named as
his account of them, since the papers are unread here. The errata page also corrects an equation it
numbers (10) to a weak form with the signs
$-\iint u\cdot\partial_t\theta - \sum\iint u_iu_j\,\partial_j\theta_i = \nu\iint u\cdot\Delta\theta
+ \iint f\cdot\theta + \iint p\,\operatorname{div}\theta$; in the text as fetched (10) is the
periodicity of $u$ and the weak form is (12), so the errata refer to an earlier numbering. Recorded as
read.

\section{What they knew, what they wanted, what we know}

Three lists, kept apart. The first two are Fefferman's statement, read; the third is the two files,
run.

\textbf{What they knew}, as the statement records it, in his order and his attributions. In two
dimensions the analogs of (A) and (B) have been known for a long time, Ladyzhenskaya, and also for
the harder case of Euler; the three-dimensional difficulties are absent there. In three dimensions (A)
and (B) hold when $u^\circ$ satisfies a smallness condition, and for general $u^\circ$ they hold,
also for $\nu = 0$, on a short interval $[0,T)$ with $T$ depending on the datum; the largest such $T$
is the blowup time, and either (A) and (B) hold or some smooth divergence-free datum has a solution
with finite blowup time, near which the velocity is unbounded. For Euler with finite blowup time the
vorticity satisfies the Beale--Kato--Majda divergence, so it blows up rapidly; many numerical
computations appear to show Euler blowup, and the extreme numerical instability of the equations
makes reliable conclusions hard to draw; he points to Bertozzi and Majda's book. Leray, 1934: weak
solutions always exist in three dimensions with suitable growth; their uniqueness for Navier--Stokes
is not known, and for Euler it is false, by Scheffer and then Shnirelman, who built weak solutions
with compact support in spacetime, a fluid at rest that starts moving with no stimulus and returns to
rest. Scheffer, then Caffarelli, Kohn and Nirenberg, then Lin: partial regularity, the singular set of
a suitable weak solution having one-dimensional parabolic Hausdorff measure zero, so that it contains
no spacetime curve; he calls it the best partial regularity theorem known so far and says it appears
very hard to go further. His closing sentence: standard methods from PDE appear inadequate, and some
deep, new ideas are probably needed. Past the statement, and as the previous chapter recorded on
2026-09-11: a 2026 paper and a Lean repository were published against (C) and (D); that chapter read
the formalized statement against his six points and found each matched, built nothing, and read
neither proof. Nothing here adds to that.

\textbf{What they wanted}: a proof of one of (A), (B), (C), (D), in the sets defined above. He says the
leeway is deliberate, to give solvers room while retaining the heart of the problem: whether smooth,
physically reasonable solutions exist for every reasonable datum, or whether some datum breaks down.
Physically reasonable is defined, not left open: smooth, with bounded energy on $\R^3$, or smooth and
periodic on the torus with the periodic pressure of the errata. The Euler case is wanted too, in his
words, and is not on the Clay list.

\textbf{What we know}, each line with its set and its scope, and nothing past the scope.

\begin{itemize}
\item On $R_{\mathrm{div}}$, a subset of $D_8$ defined in the next section, the residual of (1) with
  (2) is computable exactly, and for the Arnold--Beltrami--Childress datum it is zero at every Taylor
  order checked, orders 0 to 4, by two routes that agree to the integer. The closed form is a solution
  for every $t$ by the two identities verified; that is an instance of (B) for one datum, known since
  Arnold and Childress.
\item For a generic datum in $R_{\mathrm{div}}$ the Taylor coefficients to order 4 are exact,
  divergence-free, real-valued, and the vorticity route agrees at every order; their mode support
  grows by one in $|k|_1$ per order, measured, so no fixed horizon holds the solution. What happens
  past order 4, whether the series converges, and for how long, is not known here.
\item At $t = 0$, on both data, the energy identity holds exactly, with the nonlinear term and the
  pressure moving no energy. That is one instant; (7) is every instant.
\item The time and space scalings hold exactly on the coefficients to order 3; the viscosity is a
  parameter a scaling moves.
\item On one datum each, the solution sets of Euler and of Navier--Stokes do not contain each other's
  member, and the forced and unforced sets do not; the residuals are exact and proportional to $\nu$
  or equal to $f$. Two instances, no theorem.
\item $R_{\mathrm{div}}$ is countable and $D_8$ is not, so the exactly nameable data are a
  measure-zero island in the data class; everything the ring does is on that island.
\item The boundary kind of each object, read by probe: none on the ABC coefficients and on the exact
  ring, completeness on a generic datum's horizon, format on any decimal report of a coefficient with
  $\pi$ in it, measurement on a deposited amplitude and none on the ratio that cancels it. Through the
  chain from the datum to $u_3$ the ring carries no format kind, measurement is carried to every
  order, exactly on the ABC family, and completeness is carried and grows.
\item The disjoint-translate constructor splits the bilinear term exactly when supports are disjoint
  and not when they overlap, in one variable on a lattice of eighths. Its algebra, not their
  corollary.
\item We do not know whether (A), (B), (C) or (D) holds. We have not read the 2026 proof. We have
  computed no blowup, no weak solution, no singular set, and no solution on $\R^3$. The two files
  reach one countable island inside $D_8$ and read it exactly; they reach nothing past it.
\end{itemize}

\section{Our set, and where it sits in theirs}

The arithmetic is the engine's alone: Python integers as the bignum, decimal inputs read by
\texttt{representation.exact} as (numerator, places) pairs, and booleans; no float, no fraction
library. A field on $\R^3/\Z^3$ is a finite sum of modes $e^{2\pi i k\cdot x}$, $k\in\Z^3$, each mode
carrying a Laurent polynomial in $\pi$ with Gaussian integer coefficients, and the whole field
carrying one positive integer denominator, the way the engine's exact reader carries one count of
decimal places. The denominator is widened past a power of ten for one reason: Leray's projection
divides by $|k|^2$, and $1/3$ has no exact decimal at any scale, and a decimal scale would then have
to refuse where an integer denominator carries the value exactly. Every field is normalized by the
common divisor of its integers, so equality is equality of integers, and $\pi$ is a symbol:
Lindemann's theorem makes term-by-term zero the exact zero test. Call the ring $R$, real-valued when
the coefficient at $-k$ is the conjugate of the one at $k$, and $R_{\mathrm{div}}$ its
divergence-free vector fields. The solution is carried as Taylor coefficients at $t=0$,
\[
u_{m+1} \;=\; \nu\,\Delta u_m \;-\; P\sum_{j=0}^{m}\binom{m}{j}\,(u_j\cdot\nabla)\,u_{m-j},
\]
with $P$ Leray's projection and the pressure from the same sum, which is (1) differentiated $m$ times
at $t=0$ with (2) enforced by $P$. Nothing is rounded.

\begin{itemize}
\item $R_{\mathrm{div}} \subset D_8$: every element is a trigonometric polynomial, so smooth and
  periodic, and the divergence is an exact identity. That inclusion is what the whole of the next
  section rests on.
\item $R_{\mathrm{div}}$ meets $D_4$ only at zero: a nonzero periodic field does not decay. So the
  ring reaches (B) and (D)'s data and none of (A) and (C)'s.
\item $R_{\mathrm{div}}$ is closed under the constructors of the recurrence, products, derivatives,
  and Leray's projection with the denominator widened, so every Taylor coefficient of a solution with
  datum in $R_{\mathrm{div}}$ is in $R_{\mathrm{div}}$; shown to order 4 and true by construction at
  every order.
\item The solution itself, $u(t)$ for $t>0$, is not in $R$: for the ABC datum it is
  $e^{-4\nu\pi^2 t}u_0$, whose amplitude is not a Laurent polynomial in $\pi$ with integer
  coefficients, and for a generic datum the horizon table shows the mode set has no finite bound. $R$
  names the datum and the coefficients at one instant; $S_{1011}$ is where the solution lives, and the
  ring reaches into it only through the coefficients.
\item For the ABC datum $u_0$, $S_{1011}(u_0,0)$ is not empty, by the two identities the run verifies
  exactly, $\operatorname{curl}u_0 = 2\pi u_0$ and $(u_0\cdot\nabla)u_0 = \nabla(|u_0|^2/2)$, which
  make the closed form a solution at every $t$ since time enters only as a scalar factor. One datum's
  instance of (B), and not (B).
\end{itemize}

\section{The runs}

\texttt{exact\_navier\_stokes\_on\_torus.py}, run and exit 0.

\textbf{Positive control.} The Arnold--Beltrami--Childress field
$u_0 = (A\sin 2\pi z + C\cos 2\pi y,\; B\sin 2\pi x + A\cos 2\pi z,\; C\sin 2\pi y + B\cos 2\pi x)$
with $A,B,C = 1,2,3$ and $\nu = 1/10$, orders 0 to 4. Beltrami, $\operatorname{curl}u_0 = 2\pi u_0$,
exact. $(u_0\cdot\nabla)u_0 = \nabla(|u_0|^2/2)$, exact, and its Leray part is zero. The recurrence
returns $u_m = (-4\nu\pi^2)^m u_0$ and $p_m = (-8\nu\pi^2)^m(-|u_0|^2/2)$ at every order, to the
integer. The residual of (1) with (2) is zero at every order.

\textbf{Two routes.} The vorticity recurrence, Helmholtz's equation with velocities recovered by
Biot--Savart and no pressure anywhere, meets the velocity route at $\operatorname{curl}u_m = \omega_m$
at every order, on the exact solution and on a generic datum, and Biot--Savart returns $u_m$ from
$\omega_m$. Drawn nulls: dropping the vortex-stretching term, the two-dimensional form of the
equation, breaks the agreement at order 1 in three dimensions; flipping the pressure sign leaves a
nonzero residual.

\textbf{The energy identity} at $t=0$, $\frac{d}{dt}\int|u|^2 = -2\nu\int|\nabla u|^2$, holds
exactly on both data: on the ABC datum both sides are $-56\pi^2/5$ with $\int|u_0|^2 = 14$. On the
generic datum the two pieces that vanish are shown to vanish, $\int u\cdot(u\cdot\nabla)u = 0$ and
$\int u\cdot\nabla p = 0$, and the null is a field with divergence, $u_0 + \nabla\cos 2\pi x$, which
moves energy through a pressure. This is the mechanism behind (7) for $f = 0$, shown at one instant on
one instance; it is not (7).

\textbf{The horizon}, on a generic divergence-free datum
$g_0 = (\sin 2\pi y + \sin 2\pi z,\; \sin 2\pi z,\; \sin 2\pi x)$, not Beltrami, with a surviving
Leray part and a nonzero pressure from order 1:

\begin{center}
\begin{tabular}{@{}lrrrrr@{}}
\toprule
order & 0 & 1 & 2 & 3 & 4 \\
\midrule
largest $|k|_1$ & 1 & 2 & 3 & 4 & 5 \\
largest $|k|_\infty$ & 1 & 1 & 2 & 3 & 4 \\
modes held & 6 & 18 & 42 & 88 & 170 \\
outside $|k|_\infty \le 1$ & 0 & 0 & 16 & 62 & 144 \\
outside $|k|_\infty \le 2$ & 0 & 0 & 0 & 16 & 66 \\
\bottomrule
\end{tabular}
\end{center}

$|k|_1$ climbs by exactly one per order, because a product of modes adds their index vectors. A
truncation at any fixed radius misses some order, and the count it misses is measured, not bounded.
Every coefficient stays divergence-free and real-valued.

\textbf{Symmetries}, exact on the coefficients. Time scaling $v(x,t) = \mu u(x,\mu t)$ solves
(1)--(3) at viscosity $\mu\nu$ with $v_m = \mu^{m+1}u_m$, checked at $\mu = 3$; the wrong exponent
$\mu^m$ is refused. Space scaling $v(x,t) = \lambda u(\lambda x, \lambda^2 t)$ at the same $\nu$
keeps the period and gives $v_m = \lambda^{1+2m}u_m(\lambda x)$, checked at $\lambda = 2$. The
viscosity is a parameter a scaling moves.

\textbf{The two removals}, on instances. At $\nu = 0$ the ABC field is stationary and solves Euler
exactly; in (1) at $\nu = 1/10$ its residual is $4\nu\pi^2u_0$, nonzero; the decaying viscous field's
residual in Euler at $t = 0$ is $-4\nu\pi^2u_0$, nonzero. The field $e^{-t}g_0$ solves the forced
equation with $f$ equal to its own residual, real and periodic, nonzero at orders 0 to 2, and in the
unforced equation its residual is that $f$. Neither solution set inherits the other's membership on
these instances. This is the previous chapter's argument from the shapes of the statements, that
removing the force and removing the viscosity are different removals, shown exactly on one field
each. It says nothing about blowup and nothing about the 2026 paper's construction, which is not
represented here.

\textbf{The island.} $F_b = \sum_n 2b_n n^{-n}\cos(2\pi n x)\,e_y$, for a bit sequence $b$, is
divergence-free and real-valued at every truncation, distinct sequences give distinct fields, and the
bits read back from the coefficients. The full sums are smooth, since $n^{-n}$ beats every power of
$n$, and periodic, so they sit in $D_8$. The bit sequences are uncountable, by Cantor's diagonal as
the engine's set-theory posit reproduces it, and the ring is countable.

\section{The boundary function, asked to define itself}

\texttt{proof\_boundary\_inheritance.py}, run and exit 0. The engine's precision document names
three kinds of boundary on an exact representation, from a survey: a format boundary, which a larger
format raises; a measurement boundary, fixed upstream by whoever measured; and a completeness
boundary, which only more computation crosses. A survey assigns the kind by judgment. Here the kind is
not assigned. Two probes are applied to the object, raise our scale by one and raise our horizon by
one, and the object answers: format is what the scale probe moves, completeness what the horizon
probe moves, measurement a gap to the target that neither probe moves, none no gap. Measurement means
exactly unmoved by both of our probes, and from this side external means the same. The verdict is
only as good as the probes: the drawn null cuts the horizon probe from a completeness object and the
classifier reports measurement, the limit stated and not hidden.

Positive controls, kind fixed by construction: the number theoretic transform length on the prime
$119\cdot2^{23}+1$ at horizon $2^{25}$ reads format; a 3-ary tree count at horizon 7 of 12 reads
completeness; a deposit of six places against its true value reads measurement; $22/7$ against itself
reads none.

Their sets, read by the probes, on $u_2$ of the generic datum, whose modes reach $|k|_\infty = 2$:

\begin{center}
\begin{tabular}{@{}p{0.46\textwidth}ccc p{0.22\textwidth}@{}}
\toprule
object & gap & scale probe & horizon probe & kind \\
\midrule
$u_2$ in decimals, 30 places, modes $|k|_\infty\le1$ & yes & moves & moves & format and completeness \\
$u_2$ in decimals, 30 places, modes $|k|_\infty\le2$ & yes & moves & silent & format \\
$u_2$ in the ring, modes $|k|_\infty\le1$ & yes & silent & moves & completeness \\
$u_2$ in the ring, modes $|k|_\infty\le2$ & no & silent & silent & none \\
$u_2$ from $A$ deposited to 6 places, against true $A$ & yes & silent & silent & measurement \\
$u_1/u_0 = -4\nu\pi^2$, deposit against true $A$ & no & silent & silent & none \\
\bottomrule
\end{tabular}
\end{center}

The format kind is a property of the report, not of the object: the same coefficient reads format in
decimals, because a coefficient carrying $\pi$ has no last digit, and reads none in the ring, which
has no scale. The measurement kind is the deposit's, and it is canceled only by the ratio in which
the amplitude cancels.

Each of Fefferman's sets against the boundary function:

\begin{center}
\begin{tabular}{@{}p{0.30\textwidth}p{0.40\textwidth}p{0.26\textwidth}@{}}
\toprule
his set & what the ring does with it & kind, by probe \\
\midrule
(1), the equation behind $S_{67}$ and $S_{1011}$ & residual exact on any element of $R_{\mathrm{div}}$; solution as Taylor coefficients & none on the ABC family; completeness on a generic datum \\[0.4em]
(2), the condition on every $D$ and $S$ & exact identity, carried by every coefficient & none \\[0.4em]
(3), $u^\circ$ in $D_8$ & any element of $R_{\mathrm{div}}$; a deposited amplitude & none; measurement for the deposit \\[0.4em]
$D_4$, $F_5$, $S_{67}$ & not represented: $R_{\mathrm{div}}$ meets $D_4$ at zero & not run \\[0.4em]
(7) on the unit cell & Parseval's sum, exact; the energy identity at $t=0$, exact & none at the instant checked \\[0.4em]
$D_8$ with the errata's periodic pressure & by construction: datum, force and pressure all periodic & none \\[0.4em]
$F_{89}$, decay of $f$ in $t$ & not probed: time is carried at $t=0$ only & not run \\[0.4em]
$S_{1011}$ & by construction for every coefficient; the solution itself is outside $R$ & none on the coefficients; the solution reached only through them \\[0.4em]
$\nu > 0$ & moved exactly by the time scaling & format, in the sense a symmetry moves it \\[0.4em]
Euler, $\nu = 0$ & the ABC field is a stationary exact solution & none; nothing about blowup or the Beale--Kato--Majda integral \\[0.4em]
(A), (B) & the ring reaches one exact family in $D_8$ and finitely many orders of a generic datum & completeness; claims nothing \\[0.4em]
(C), (D) & the 2026 paper's datum and force are not represented, and the paper is unread past what the previous chapter read & not run \\[0.4em]
the weak solutions and the singular set & not represented & not run \\
\bottomrule
\end{tabular}
\end{center}

\section{Constructors, and what inherits their proof}

A ring element is built by a fixed set of constructors: exact integer arithmetic with one
denominator per field, the mode convolution, a linear convolution on $\Z^3$ that never aliases, the
derivative multiplier $2\pi i k_j$, Leray's projection, and the Taylor recurrence. A coefficient
built only from proven constructors inherits their proof. What each boundary kind does through the
chain from the datum to $u_3$, measured:

\begin{center}
\begin{tabular}{@{}p{0.20\textwidth}p{0.22\textwidth}p{0.30\textwidth}p{0.24\textwidth}@{}}
\toprule
kind & at the datum & through the chain & inherited? \\
\midrule
format, in the ring & absent & scale probe silent at every order & no format kind exists to inherit \\[0.4em]
format, in decimals & absent: the datum's coefficients are rational & moves from order 1 on: the derivative brings $\pi$ in & introduced by the constructors, then carried \\[0.4em]
measurement, ABC & the deposit's gap & $\mathrm{gap}_m = (-4\nu\pi^2)^m\,\mathrm{gap}_0$, exactly & carried exactly, unamplified; canceled by the ratio \\[0.4em]
measurement, generic & the deposit's gap & nonzero at every order & carried; no order lowers it \\[0.4em]
completeness & 0 modes outside $|k|_\infty\le1$ & 0, 0, 16, 62 & carried and growing \\
\bottomrule
\end{tabular}
\end{center}

Two more inheritances, one run and one read.

\textbf{The disjoint-translate constructor.} The previous chapter's account of the 2026 paper's
Corollary 10.6 is that the periodic case is built by summing integer translates of a compactly
supported solution whose supports stay disjoint. What that sum inherits from its pieces through the
bilinear term rests on one exact fact, and the posit shows it on piecewise polynomials with integer
coefficients on a lattice of eighths, $y = 8x$: with pieces $(y-a)^2(b-y)^2$ on $[0,2]$ and $[4,6]$,
the bilinear term of the sum splits into the pieces' bilinear terms by coefficients and by the exact
integral of the squared difference, which is $0$; with the second piece moved to $[1,3]$, the drawn
null, both routes refuse and the integral of the squared cross term over $y$ is $156928/5005$. The
linear terms split regardless. The construction is theirs. Only its algebra is run here, in one
variable, and nothing here checks the paper's own proof of that corollary, which was not read.

\textbf{Fefferman's statement to the Lean statement.} That is a reading, not a computation, and the
previous chapter did it on six points and found each matched. Nothing here adds to it, and nothing
here was machine-checked either.

\section{The figure}

The figure below is drawn by the engine, not by hand:
\texttt{maint/\allowbreak{}texbuild/\allowbreak{}navier\_stokes\_sets\_figure.py} in the
\texttt{anchor\_sift} tree writes the PDF itself, with no drawing library, and takes the right panel's
numbers from running the torus example at that moment. The left panel is the shape of the four
alternatives read off the set definitions above: the base is the plane of pairs (datum, force), the
front edge is force zero, a stalk over a point is the solution set for that pair. (A) and (B) speak of
the edge, (C) and (D) of the plane, and the edge lies in the plane. Every stalk but one is drawn as a
question, because that is what it is; the one drawn solid is the ABC datum, one instance.

\IfFileExists{../../build/theory/figures/navier_stokes_sets.pdf}{%
\begin{figure}[htbp]
\centering
\includegraphics[width=\linewidth]{../../build/theory/figures/navier_stokes_sets.pdf}
\caption{Left: the plane the breakdown alternatives speak of holds the edge the existence
alternatives speak of; every stalk but one is a question. Right: modes held by the Taylor
coefficients $u_0$ to $u_4$ of a generic datum, from the run, with the count outside two fixed boxes
inside each bar. Generated by \texttt{navier\_stokes\_sets\_figure.py}; it claims nothing about (A),
(B), (C) or (D).}
\end{figure}}{%
The figure is absent from this build. Generate it with
\texttt{python maint/texbuild/navier\_stokes\_sets\_figure.py} from the \texttt{anchor\_sift} root,
which writes \texttt{build/theory/figures/navier\_stokes\_sets.pdf}, and build again.}

\section{Prior art, named with respect}

Every object above is the field's. The problem statement, its numbered conditions, the four
alternatives and the errata are Charles Fefferman's for the Clay Mathematics Institute, read in full.
The equations are Navier's and Stokes's, the inviscid case Euler's. The field with
$\operatorname{curl}u$ proportional to $u$ is Beltrami's; the three-term example is Arnold's (1965)
and Childress's (1970), reported from memory of the literature and unread here, so the file verifies
the solution itself and does not rest on the citation. The projection onto divergence-free fields and
the pressure it defines are Leray's. The recovery of a velocity from its vorticity is the Biot--Savart
law; the vorticity equation is Helmholtz's. The energy sum over modes is Parseval's. The
transcendence of $\pi$ is Lindemann's (1882). The Taylor recurrence is the classical power-series
method of Cauchy and Kovalevskaya. The Beale--Kato--Majda criterion is quoted from Fefferman's
statement and not used. Cantor's diagonal and the Moore closure are credited in the engine's
set-theory posit. Machin's formula (1706) supplies $\pi$ where a decimal is needed. The 2026 paper on
finite-time blowup and its Lean repository are described only as the previous chapter described them,
and that chapter states what it did not read. Result checking by an independent route is Blum, Luby
and Rubinfeld's, credited in the engine's analytic number theory workbook. Nothing in either file is
new mathematics, and neither file wants it to be.

\section{Withdrawn}

Four claims were withdrawn in the course of this chapter's work, and each stays on the page with what
killed it.

\textbf{Withdrawn.} The first build of both files, which carried the coefficients as Gaussian
rationals from a fraction library and the binomials from a math library. \textbf{What killed it:}
Douglas Quigg, on reading it. The engine's arithmetic is Python integers as the bignum, decimal text
read by the exact reader, and booleans, with no other arithmetic beside them; a library that carries
rationals for us is a second arithmetic under the one the tree is built to trust, and the exact path
exists so that there is no second arithmetic. Rebuilt on integers with one denominator per field, and
every result above was reproduced to the integer. The one number that changed is the overlap
integral, $613/44024445576151040$ over $x$ in the first build and $156928/5005$ over $y = 8x$ in this
one, because the lattice moved to eighths; both are nonzero, and that is all the null asks of them.
Recorded first because it was the largest fault.

\textbf{Withdrawn.} That the mode radius $|k|_\infty$ of the generic datum's Taylor coefficients grows
at every order. \textbf{What killed it:} the run. At order 1 the radius is 1, the same as at order 0,
because the first products of unit modes land on modes like $(0,1,1)$, whose sup-norm is still 1.
Replaced by the measured statement: $|k|_1$ climbs by exactly one per order, and $|k|_\infty$ never
falls and ends higher. The first wording was derived from the shape of a convolution and not drawn
from the output.

\textbf{Withdrawn.} That the format kind, in a decimal report, is present at every order of the
chain. \textbf{What killed it:} the run. At the datum the coefficients are $\pm i/2$, terminating
rationals, and the scale probe is silent there. The kind appears at order 1, when the derivative
multiplies by $2\pi i k$. Replaced by the measured statement that the constructors introduce it. The
better finding was in the refusal.

\textbf{Withdrawn}, for the duration of one run. That two bit patterns produced two distinct fields
whose bits read back. The first build of the island summed $\cos 2\pi x$ for every $n$ in place of
$\cos 2\pi n x$, so every bit landed on one mode. The check that reads the bits back refused it; the
harmonic was corrected and the check passed. Recorded because a wrong construction had produced a
true reading of distinctness by accident, and only the second check caught it.

\section{What was not verified, and what is still open}

\textbf{Nothing here is a theorem about the problem.} Every statement in the runs section is an
exact computation on one or two data to a stated order, and the boundary and inheritance tables are
about those computations. The whole-space conditions (4) to (7) need a representation with decay on
$\R^3$, which the periodic ring is not; a compactly supported piecewise-polynomial representation in
three variables would reach the disjoint-translate constructor in its own setting, and the
one-variable algebra above has the shape of that setting and is not it. Condition (9), decay of the
force in time, needs the time dependence carried past $t = 0$, which Taylor coefficients at one
instant do not see. The horizon count is measured to order 4 on one datum; whether $|k|_1$ climbing by
one per order is the general rule it looks like or the property of this datum is stated above as
measured, on this datum, to this order.

\textbf{The 2026 proof was not read} past what the previous chapter read, and nothing here rests on
it. \textbf{Arnold and Childress were not read}, and the ABC solution is verified in the file instead
of cited for. \textbf{Nothing was machine-checked.}
